\documentclass[10pt,a4paper]{amsart}
\usepackage[margin=19mm]{geometry}
\usepackage{amsmath,amssymb,mathtools}
\usepackage[scr=rsfs,cal=euler]{mathalfa}
\usepackage{xfrac}
\usepackage[unicode,hidelinks,bookmarks=false]{hyperref}
\newcommand{\ie}{i.e.\ }
\newcommand{\cf}{cf.\ }
\newcommand{\resp}{respectively\ }
\newcommand{\Subsection}{\subsection}
\providecommand{\nobreakdash}{\nobreak\hbox{-}}
\setlength{\emergencystretch}{2em}
\multlinegap0pt
\usepackage{mathtools}
\usepackage{amssymb}
\newtheorem{lem}{Lemma}
\newcommand{\R}{\mathbb{R}}                 % real numbers
\newcommand{\col}{\colon \thinspace}        % colon : with small space for maps
\title{Deformations of harmonic maps with conical singularities}
\author{Dominik Gutwein, Thibault Langlais}
\date{Source: arXiv:2609.20588v1 (CC BY 4.0)}
\begin{document}
\maketitle

\noindent\emph{Source and licence.} Deformations of harmonic maps with conical singularities. Version arXiv:2609.20588v1 (CC BY 4.0), distributed by arXiv under the Creative Commons Attribution 4.0 International licence, http://creativecommons.org/licenses/by/4.0/, as recorded in the arXiv licence metadata for this exact version and shown on the abstract page https://arxiv.org/abs/2609.20588v1. Full licence notice: \texttt{licenses/arxiv-2609.20588-CC-BY-4.0.txt}.\par\smallskip

\noindent\textit{Editorial scope.} Counted unit: the article's lem values (source0.tex lines 2991--3021). The article's complete original proof of it is delivered with it (source0.tex lines 3023--3146, one \texttt{proof} environment of 124 lines). No mathematics is written, repaired or re-derived: the statement and the proof are the article's own lines, in the article's own notation, and no retained line is substituted or re-flowed.  No further source-local context is needed: every cross-reference of every retained line resolves inside the two delivered spans.  Nothing was omitted from any retained span, so the package carries no omission record.  No retained line contains a citation, so the wrapper carries no bibliography and no bibliography entry.  No retained line contains a figure environment, an image-inclusion command or drawing code, so no image is delivered and the package declares no asset dependency.  Editorial formatting only, in addition: an AMS \texttt{amsart} preamble that restates the article's own declarations and macros used by the retained lines, namely the environments \texttt{lem} on separate counters, together with the standard packages those lines need.  The retained environments print in this document as Lemma 1 in order of appearance, whereas the article prints the same items with the numbering of its own full document.  The complete native source of the article as distributed in the arXiv e-print for this exact version is bundled unchanged as \texttt{sources/210887/source0.tex}.
\begin{lem}     \label{lem: Point (ii) of admissible values}
    Let us fix $v > 0$, $\sigma \in \R \backslash \{0\}$ and a smooth function $\psi(\varepsilon,t) \col [0,1) \times [0,\infty) \rightarrow \R$ with all derivatives uniformly bounded and such that 
    \begin{equation*}
        \psi(0,t) = \frac{\partial \psi}{\partial \varepsilon}(0,t) = 0 .
    \end{equation*}
    Consider $X(\varepsilon,\cdot), Y(\varepsilon,\cdot) \col [0,\infty) \rightarrow \R^2$, the unique solutions of the ODE
    \begin{equation*}
        - U^{\prime\prime}(t) - \psi(\varepsilon,t) U^\prime(t) + \begin{pmatrix} 1 & \sigma \varepsilon \chi(t) \\ \sigma \varepsilon \chi(t) & 0 \end{pmatrix} U(t) = 0
    \end{equation*}
    satisfying the initial conditions
    \begin{equation*}
        X_\pm(\varepsilon,1) = 
        \begin{pmatrix}
            1 \\
            0
        \end{pmatrix}, ~ X^\prime_\pm(\varepsilon,1) = 
        \begin{pmatrix}
            v \\
            0
        \end{pmatrix}, ~~~ \text{and} ~~~ Y_\pm(\varepsilon,1) = 
        \begin{pmatrix}
            0\\
            1
        \end{pmatrix}, ~ Y^\prime_\pm(\varepsilon,1) = 
        \begin{pmatrix}
            0\\
            0
        \end{pmatrix}.
    \end{equation*}
    Then if $T \geq 0$ is large enough, there exists $\varepsilon(T) > 0$ such that for all $\varepsilon \in (0,\varepsilon(T))$, $X^\prime(\varepsilon,T)$ and $Y^\prime(\varepsilon,T)$ are linearly independent in $\R^2$.
\end{lem}

\begin{proof}
    Since the resolvent of the ODE is a smooth function of $(\varepsilon,t) \in [0,1) \times [0,\infty)$, there are expansions
    \begin{align*}
        X(\varepsilon,t) & = X_0(t) + \varepsilon X_1(t) + \frac{\varepsilon^2}{2} X_2(t) + \varepsilon^3X^{(3)}(\varepsilon,t) , \\
        Y(\varepsilon,t) & = Y_0(t) + \varepsilon Y_2(t) + \frac{\varepsilon^2}{2} Y_2(t) + \varepsilon^3Y^{(3)}(\varepsilon,t)
    \end{align*}
    where $X^{(3)}(\varepsilon,t)$ and $Y^{(3)}(\varepsilon,t)$ are smooth functions of $(\varepsilon,t) \in [0,1) \times [0,\infty)$. Moreover, 
    \begin{equation*}
        X_1(0) = X_1^\prime(0) = X_2(0) = X_2^\prime(0) = \begin{pmatrix}0 \\ 0\end{pmatrix}, ~~~ Y_1(0) = Y_1^\prime(0) = Y_2(0) = Y_2^\prime(0) = \begin{pmatrix}0 \\ 0\end{pmatrix}
    \end{equation*}
    and
    \begin{equation*}
        X_0(t) = \begin{pmatrix}
            \cosh(t) + v \sinh(t) \\
            0
        \end{pmatrix}, 
        ~~
        Y_0(t) = \begin{pmatrix}
            0 \\ 1
        \end{pmatrix} \cdot
    \end{equation*}
    The next orders of the expansion may be determined as solutions of an ODE, by differentiating the original ODE with respect to $\varepsilon$. For $X_1(t)$, we obtain the ODE
    \begin{equation*}
        - X_1^{\prime\prime}(t) + \begin{pmatrix} 1 & 0 \\ 0 & 0\end{pmatrix}X_1(t) + \sigma \chi(t) \begin{pmatrix} 0 & 1 \\ 1 & 0\end{pmatrix} X_0(t) = 0 .
    \end{equation*}
    The unique solution of the ODE, with the correct initial conditions, is
    \begin{equation*}
        X_1(t) = \begin{pmatrix}
            0 \\ \sigma \Psi_v(t)
        \end{pmatrix}, ~~~ \Psi_v(t) = \int_0^t \int_0^s \chi(u)(\cosh(u) + v \sinh(u)) \mathrm{d}u\mathrm{d}s .
    \end{equation*}
    Similarly, for $Y_1$ we have the equation 
    \begin{equation*}
        - Y_1^{\prime\prime}(t) + \begin{pmatrix} 1 & 0 \\ 0 & 0\end{pmatrix}Y_1(t) + \sigma \chi(t) \begin{pmatrix} 0 & 1 \\ 1 & 0\end{pmatrix} Y_0(t) = 0 .
    \end{equation*}
    If we denote by $\Phi(t)$ the unique solution of the ODE
    \begin{equation*}
        f^{\prime\prime}(t) = f(t) + \chi(t), ~~~ f(0) = f^\prime(0) = 0,
    \end{equation*}
    then we have
    \begin{equation*}
        Y_1(t) = \begin{pmatrix}
            \sigma \Phi(t) \\ 0
        \end{pmatrix} \cdot
    \end{equation*}
    For our purpose, we only need to determine $Y_2(t)$, as $X_2(t)$ will not be relevant. It is determined by the ODE
    \begin{equation*}
        -Y_2^{\prime\prime}(t) + \frac{\partial^2\psi}{\partial\varepsilon^2}(0,t) Y_0^\prime(t) + \begin{pmatrix} 1 & 0 \\ 0 & 0 \end{pmatrix} Y_2(t) + 2 \sigma \chi(t) \begin{pmatrix} 0 & 1 \\ 1 & 0 \end{pmatrix} Y_1(t) = 0 .
    \end{equation*}
    Since $Y_0^\prime(t) = 0$ for all $t$, we can see that
    \begin{equation*}
        Y_2(t) = \begin{pmatrix}
            0 \\
            2 \sigma^2 \Xi(t)
        \end{pmatrix}, ~~~ \Xi(t) = \int_0^t \int_0^s \chi(u)\Phi(u)  \mathrm{d}u\mathrm{d}s .
    \end{equation*}

    Gathering the previous results, we obtain
    \begin{equation*}
        X(\varepsilon,t) = \begin{pmatrix}
            f_v(t) + \mathcal{O}(\varepsilon^2) \\ \varepsilon \sigma \Psi_v(t) + \mathcal{O}(\varepsilon^2)
        \end{pmatrix}, ~~~ f_v(t) = \cosh(t) + v \sinh(t)
    \end{equation*}
    and 
    \begin{equation*}
        Y(\varepsilon,t) = \begin{pmatrix}
            \varepsilon \sigma \Phi(t) + \mathcal{O}(\varepsilon^3) \\ 1 +\varepsilon^2 \sigma^2 \Xi(t) + \mathcal{O}(\varepsilon^3)
        \end{pmatrix} \cdot
    \end{equation*}
    Now we can compute
    \begin{align*}
        \det \begin{pmatrix} X^\prime(\varepsilon,t) & Y^\prime(\varepsilon,t) \end{pmatrix} & = \det \begin{pmatrix}
            f^\prime_v(t) + \mathcal{O}(\varepsilon^2) & \varepsilon \sigma \Phi^\prime(t) + \mathcal{O}(\varepsilon^2) \\ \varepsilon \sigma \Psi_v^\prime(t) + \mathcal{O}(\varepsilon^2) & \varepsilon^2 \sigma^2 \Xi^\prime(t) + \mathcal{O}(\varepsilon^3) 
        \end{pmatrix} \\
        & = \varepsilon^2 \sigma^2 (f^\prime_v(t) \Xi^\prime(t)-\Psi^\prime_v(t) \Phi^\prime(t)) + \mathcal{O}(\varepsilon^3) .
    \end{align*}
    Therefore, it is enough to prove that the quantity 
    \begin{equation*}
        D_v(t) \coloneqq f^\prime_v(t) \Xi^\prime(t)-\Psi^\prime_v(t) \Phi^\prime(t)
    \end{equation*}
    does not vanish for $t$ large enough. 

    First, let us remark that for $t \geq 1$
    \begin{align*}
        \Psi^\prime_v(t) & = \int_0^t \chi(s) (\cosh(s) + v \sinh(s)) \mathrm{d}s = \sinh(t) + v \cosh(t) + C_v = f^\prime_v(t) + C_v
    \end{align*}
    where 
    \begin{equation*}
        C_v = - \sinh(1) - v \cosh(1) + \int_0^1 \chi(s) (\cosh(s) + v \sinh(s)) \mathrm{d}s .
    \end{equation*}
    Therefore,
    \begin{equation*}
        D_v(t) = f^\prime_v(t) (\Xi^\prime(t)-\Phi^\prime(t)) + C_v \Phi^\prime(t) .
    \end{equation*}
    On the other hand, the quantity $\Xi^\prime(t) - \Phi^\prime(t)$ satisfies the ODE
    \begin{equation*}
        (\Xi^\prime(t)-\Phi^\prime(t))^\prime = \Xi^{\prime\prime}(t) - \Phi^{\prime\prime}(t) = \chi(t) \Phi(t) - \Phi^{\prime\prime}(t)
    \end{equation*}
    and by definition $\Phi^{\prime\prime}(t) = \Phi(t) + \chi(t)$, whence
    \begin{equation*}
        (\Xi^\prime(t)-\Phi^\prime(t))^\prime = (\chi(t)-1)\Phi(t) - \chi(t) .
    \end{equation*}
    For $t \geq 1$, $\chi(t) = 1$ and therefore for $t \geq 1$ we have
    \begin{equation*}
        \Xi^\prime(t)-\Phi^\prime(t) = A - t , ~~~ A = 1 + \Xi^\prime(1) - \Phi^\prime(1) .
    \end{equation*}
    It follows that for $t \geq 1$, 
    \begin{equation*}
        D_v(t) = (A- t) f^\prime_v(t) + C_v \Phi(t) .
    \end{equation*}
    As $t \rightarrow \infty$, $f^\prime_v(t) = \sinh(t) + v \cosh(t) = \frac{1+v}{2} \mathrm{e}^t + \mathcal{O}(\mathrm{e}^{-t})$. On the other hand, since 
    \begin{equation*}
        \Phi^{\prime\prime}(t) = \Phi(t) + \chi(t) \leq \Phi(t) + 1
    \end{equation*}
    and the unique solution of the ODE
    \begin{equation*}
        f^{\prime\prime}(t) = f(t) + 1, ~~~ f(0) = f^\prime(0) = 0
    \end{equation*}
    is $f(t) =\cosh(t)-1$, we have $0 \leq \Phi(t) \leq \cosh(t)-1 = \frac{1}{2}\mathrm{e}^t + \mathcal{O}(\mathrm{e}^{-t})$, and hence
    \begin{equation*}
        D_v(t) = (A-t) \frac{1+v}{2} \mathrm{e}^t + \mathcal{O}(\mathrm{e}^t) \sim -\frac{(1+v)t}{2} \mathrm{e}^t \rightarrow - \infty
    \end{equation*}
    as $t \rightarrow \infty$, which finishes the proof.
\end{proof}

\end{document}
